\section{PHẦN KẾT LUẬN}

\subsection{1. Đánh giá tổng quan về tính khả thi và sự đồng bộ của kế hoạch dự án}

Bản Kế hoạch Quản lý Dự án đối với đề tài \textit{``Dự án xây dựng bộ tài liệu tư vấn phân luồng và định hướng nghề nghiệp cho học sinh cuối cấp trung học cơ sở''} đã thiết lập một phương án can thiệp đồng bộ về phạm vi, tiến độ, ngân sách, chất lượng, nhân sự, rủi ro, truyền thông và mua sắm nhằm giải quyết trực tiếp nhu cầu học liệu hướng nghiệp cho học sinh lớp 9.

Mười chương nội dung của bản kế hoạch được liên kết chặt chẽ theo một mạch thông số thống nhất giữa các hợp phần quản trị:

Từ kết quả phân tích Cây vấn đề và Cây mục tiêu (Chương 1), ma trận lựa chọn chiến lược đã luận chứng định lượng lý do lựa chọn phương án biên soạn và xuất bản ấn phẩm chuẩn hóa (đạt \textbf{8,55/10 điểm}). Trên cơ sở đó, Điều lệ dự án và cơ chế kiểm soát thay đổi tổng thể (Chương 2) xác lập khung pháp lý cho Ban Quản lý dự án triển khai đề án trong 90 ngày lịch (từ ngày 02/03/2026 đến ngày 30/05/2026) với ngân sách 185.000.000 đồng.

Tuyên bố phạm vi và Cấu trúc phân rã công việc (Chương 3) chia nhỏ dự án thành 5 giai đoạn và 20 gói công việc cụ thể (\texttt{1.1} đến \texttt{5.4}) nhằm tạo ra 2.000 cuốn tài liệu khổ 17 $\times$ 24 cm (gồm 80 trang ruột và 4 trang bìa ngoài) kèm bộ tệp tin chế bản điện tử gốc. Hai mươi gói công việc này được tích hợp thành 10 công việc mạng (A đến J) tại Chương 4, qua đó xác định chính xác đường găng gồm 9 công việc có tổng độ dài 90 ngày và phát hiện khoản dự trữ 20 ngày của công việc mua sắm in ấn (D) để chủ động điều hòa nguồn lực.

Hệ thống 20 gói công việc trong Từ điển phân rã công việc tiếp tục được ánh xạ trực tiếp sang Bảng dự toán chi phí từ dưới lên tại Chương 5 (áp dụng định mức của Thông tư số 51/2019/TT-BTC) và gắn với đường mức chi phí cơ sở qua 5 đợt giải ngân theo mốc nghiệm thu. Tại Chương 6, công tác kiểm soát chất lượng được định lượng hóa bằng bài toán lấy mẫu thống kê $n = 384$ cuốn (độ tin cậy 95\%, sai số chấp nhận được 5\%) cùng bảng phân tích 5 nhóm chi phí chất lượng trị giá \textbf{41.200.000 đồng} (chiếm \textbf{22,27\%} tổng ngân sách).

Ở góc độ tổ chức thực thi, Chương 7 thiết lập cơ cấu ma trận cân bằng gồm 4 tiểu ban (với 4 nhóm tác giả chuyên môn) và ma trận phân bổ trách nhiệm bảo đảm mỗi công việc chỉ có duy nhất một cá nhân giữ quyền phê duyệt (\texttt{A}). Chương 8 định lượng hóa bài toán lựa chọn nhà thầu in bằng mô hình Cây quyết định ($\text{EMV}_1 = +16.000.000\text{ đồng} > \text{EMV}_2 = -10.000.000\text{ đồng}$), làm tiền đề trực tiếp cho Chương 10 tổ chức chào hàng cạnh tranh rút gọn gói thầu in ấn (dự toán 95.000.000 đồng, trúng thầu 94.000.000 đồng theo Luật Đấu thầu số 22/2023/QH15). Toàn bộ luồng thông tin phối hợp, báo cáo tích hợp tại mốc ngày thứ 60 và thủ tục đóng gói \textbf{03 bộ hồ sơ lưu trữ gốc} được chuẩn hóa trọn vẹn tại Chương 9.

\subsection{2. Khuyến nghị khai thác sản phẩm sau bàn giao}

Để 2.000 cuốn tài liệu và bộ tệp tin chế bản điện tử gốc phát huy hiệu quả sau khi dự án kết thúc vào ngày 30 tháng 5 năm 2026, Ban Quản lý dự án đề xuất ba nhóm khuyến nghị đối với các đơn vị tiếp nhận:

\textbf{Đối với Phòng Giáo dục và Đào tạo:} Ban hành văn bản hướng dẫn các trường trung học cơ sở sử dụng bộ tài liệu trong các tiết chủ đề hướng nghiệp thuộc môn Hoạt động trải nghiệm, hướng nghiệp lớp 9; đồng thời đăng tải phiên bản tài liệu điện tử (từ bộ tệp tin chế bản gốc đã nghiệm thu) lên Cổng thông tin điện tử của Phòng để học sinh và phụ huynh toàn địa bàn có thể tải về tra cứu miễn phí.

\textbf{Đối với Ban Giám hiệu và giáo viên năm trường trung học cơ sở thụ hưởng:} Phân bổ ngay 2.000 cuốn tài liệu về các lớp 9 và thư viện trường; hướng dẫn học sinh thực hiện bài trắc nghiệm Holland 36 câu hỏi tại Phần 1 và khuyến khích các em mang cuốn tài liệu về nhà để cùng cha mẹ điền phiếu đối thoại chọn trường tại Phần 4 trước thời điểm đăng ký nguyện vọng thi tuyển sinh vào lớp 10.

\textbf{Đối với các cơ sở giáo dục nghề nghiệp trên địa bàn:} Phối hợp định kỳ hàng năm với Phòng Giáo dục và Đào tạo để cung cấp số liệu cập nhật về chỉ tiêu tuyển sinh hệ chín cộng, ngành nghề đào tạo mới và cam kết việc làm sau tốt nghiệp, làm cơ sở cho việc cập nhật tờ phụ lục thông tin hướng nghiệp trong các năm học tiếp theo.

